%% kapa: exam=1 round=35 year=2024 period=1 form=A subject=civil_law subject_ko=민법 seq=12 number=12
%% source: Q-Net 감정평가사 기출문제 (cst00309) artlSeq=5213229 file=raw/1008/5213229/1교시 문제지 원본.pdf
%% conversion: kordoc 4.24.51 -> tools/gen1.py (auto); review=proofread-sonnet; flags=none
\begin{question}{12}
甲은 乙의 임의대리인이다. 이에 관한 설명으로 옳은 것은? (다툼이 있으면 판례에 따름)
\begin{choices}
  \item 甲이 乙로부터 매매계약체결의 대리권을 수여받아 매매계약을 체결하였더라도 특별한 사정이 없는 한 甲은 그 계약에서 정한 중도금과 잔금을 수령할 권한은 없다.
  \item 甲이 乙로부터 금전소비대차 계약을 체결할 대리권을 수여받은 경우, 특별한 사정이 없는 한 甲은 그 계약을 해제할 권한도 가진다.
  \item 乙이 사망하더라도 특별한 사정이 없는 한 甲의 대리권은 소멸하지 않는다.
  \item 미성년자인 甲이 乙로부터 매매계약체결의 대리권을 수여받아 매매계약을 체결한 경우, 乙은 甲이 체결한 매매계약을 甲이 미성년자임을 이유로 취소할 수 없다.
  \item 甲이 부득이한 사유로 丙을 복대리인으로 선임한 경우, 丙은 甲의 대리인이다.
\end{choices}
\end{question}
\answer{4}
